% Generated by scripts/build-paper.py — edit paper/src/survey.src.md
\section{Conclusion}\label{sec:conclusion}

Typed decision models make an old requirement concrete: software needs answers it can parse, probabilities it can act on and a principled route to escalate when unsure. Three weeks of evidence on Jev show that the first requirement is met by construction; the second holds on average for many tasks but must be verified per workload, and coherence across questions cannot be assumed; and the third works as a first pass with thresholds fixed in advance rather than as a replacement for stronger models. The same evidence draws a sharp boundary: typed decisions read well and compute poorly, and they are no safer than the state they are given. An open ecosystem already reproduces the interface in many ways and beats the commercial model on specialised tasks, which is a strength for research and a warning against reading any of them as the commercial model. The most useful next steps are version-pinned, same-protocol experiments that measure binding, rejection, coherence and escalation together, with raw predictions released.
